\documentclass[../../thesis.tex]{subfiles}

\graphicspath{{Figures/}}

\begin{document}
\chapter{Frequency Map-Making} \label{chap:FMM}

In Chapter \ref{chap:spectral_imaging}, we presented the frequency dependency of QUBIC's beam, and about the possibility to use this property, called \emph{spectral imaging}, to split the physical bandwidth into frequency sub-bands to increase the spectral resolution of the instrument. We now introduce a map-making algorithm based on spectral imaging, to reconstruct multiple frequency maps of the sky. We call this algorithm \emph{Frequency Map-Making} (FMM).
We start by introducing the general fundamentals of map-making for CMB experiment, followed by the description of the numerical simulation implementation for QUBIC Instrument. The resulting maps are then analysed to understand the behaviour of FMM : numerical cost, noise structure and anti-correlations introduced by spectral imaging. 

\rmk{JC, can you tell me about the exact chronology ?}
\personal{The content of this chapter is the result multiple people within QUBIC Collaboration. HERE, I WILL ADD A COMPLETE DESCRIPTION OF THE PEOPLE WHO WORKED ON MAP-MAKING FOR QUBIC, Non-exhausted list : JC Hamilton, Pierre Chanial, Mikhail Stolpovskiy, Louise Mousset (je ne vois pas vraiment de développement dans sa thèse), Mathias Regnier.
During my first year, I contributed with Mathias Regnier to the development and the implementation of the Frequency Map-Making algorithm. This work is presented in Mathias' thesis \cite{mathiasthesis}, and led to a publication \cite{FMM}. During my second year, I completly refactored and cleaned the FMM code, then I built a pipeline allowing simulation, reconstruction and analysis for FMM, easy to use from a single parameters file. During my third year, I extended the analysis of the numerical parameters and their effect on the reconstruction, leading to a forecast-ready pipeline for QUBIC instrument. My contributions to FMM are furthered developped in Chapters \ref{chap:qubicsoft}, \ref{chap:tod}, \ref{chap:xspectrum} and \ref{chap:forecast}.}


\rmk{PB: same subsection names than in Mathias' Thesis..., I am afraid that it is too close to his chapter.}
\section{Map-Making Principle} \label{section:mapmaking}

Map-making algorithms lie at the heart of the data analysis pipeline of all telescopes. These algorithms use Time-Ordered Data (TOD), datasets in which each sample corresponds to the sky signal measured by the instrument at a given time, to reconstruct maps of the sky. This reconstruction relies on the pointing direction of the telescope at each time sample, together with a model of the instrument. In this section, we develop the standard map-making formalism used by most CMB experiments and its adaptation to the specific features of QUBIC.

\subsection{Map-Making Solution}

The purpose of map-making is to reconstruct sky maps from telescope data. In CMB experiments, the data are commonly referred to as Time-Ordered Data (TOD), since the temporal dependence of the measurements is a key element in the reconstruction process through the scanning strategy of the instrument. This differs significantly from optical telescopes, which typically rely on long exposure observations to produce accurate and low-noise images of the sky. In contrast, CMB instruments continuously scan the sky following specific observing strategies. Different scanning strategies present different advantages and drawbacks, but their primary purpose is to mitigate systematic effects, such as electronic or cryogenic fluctuations, by observing the same sky location under different instrumental conditions. Scanning strategies also help reduce the impact of atmospheric $1/f$ noise, which dominates at low frequencies. In addition, modern CMB experiments aim to measure the polarisation of the CMB. This requires observing the same region of the sky with different detector orientation angles in order to distinguish the Stokes parameters $Q$ and $U$ (see \ref{subsub:polarisation}).

The most general way to describe the data acquisition at one time is : 

\begin{equation}
    d(\vec{u}(t)) = \int_{\nu_{min}}^{\nu_{max}} \mathcal{H}_\nu(\vec{u}(t)) s_\nu d\nu + n(\vec{u}(t)),
\end{equation}

where $\vec{u}(t)$ is the pointing direction on the sky at time $t$, $d$ is the TOD at time $t$, $\nu_{min}$ and $\nu_{max}$ the physical frequency bandwidth of the instrument, $\mathcal{H}_\nu$ is the acquisition operator of the instrument, which models the data acquisition of the instrument (including convolution, PSF, optics, electronics, etc), $s_\nu$ is the sky map and $n$ is the detector and photon noise contribution at time $t$. $\mathcal{H}_\nu$ is a very complex object, which will be discussed in details in (\rmk{I don't know where yet, it can be in this chapter or in qubicsoft one}). 

For simplicity, the scanning strategy is included in the acquisition operator. The entire TOD acquisition can then be written as : 

\begin{align} \label{eq:acquisition}
    \vec{d} 
    &= \int_{\nu_{min}}^{\nu_{max}} \mathcal{H}_\nu  \vec{s_\nu} d\nu + \vec{n} \notag \\
    &= H\vec{s} + \vec{n},
\end{align}

defining $H$ as the acquisition operator integrating into the frequency band. TThe goal is now to estimate the unknown sky $\vec{s}$ from the data $\vec{d}$ and the instrument model encoded in the acquisition operator $H$.
Assuming a gaussian noise $\vec{n}$, all the statistical information is stored in the noise covariance matrix $N=\left\langle\vec{n}\vec{n}^{~T}\right\rangle$. Its probability distribution, noting $N_{pointings}$ the number of time sample, is then written :

\begin{equation}
    P(\vec{n}) = \frac{1}{\left|(2\pi)^{N_{pointings}}N\right|^{1/2}} \exp\left(-\frac{1}{2}\vec{n}^{~T} N^{-1}\vec{n}\right).
\end{equation}

The best solution for map-making equation \ref{eq:acquisition} is obtained by maximising the likelihood function giving the probability of getting data $\vec{d}$ from a sky $\vec{s}$, noted $\mathcal{L} \equiv P(\vec{d} | \vec{s})$, as detailed in \cite{wright1996}. Using Bayes theorem~\cite{bayes1763} :

\begin{align}
    \mathcal{L} \equiv P(\vec{d} | \vec{s}) = \frac{P(\vec{s}|\vec{d})P(\vec{d})}{P(\vec{s})} \propto P(\vec{s}|\vec{d}).
\end{align}

As said above, we expect the noise to follow a Gaussian distribution. The probability to observe the sky according to our data is then equal to the probability distribution of the noise, meaning : $\mathcal{L} \propto P(\vec{n})$. 
Considering equation \ref{eq:acquisition}, $\vec{n} = \vec{d} - \vec{s}$. The likelihood function can then be written as : 

\begin{equation}
    \mathcal{L} \propto P(\vec{n}) = \frac{1}{\left|(2\pi)^{N_{pointings}}N\right|^{1/2}} \exp\left(-\frac{1}{2}(\vec{d} - H\vec{s})^{~T} N^{-1}(\vec{d} - H\vec{s})\right).
\end{equation}

Maximising the likelihood $\mathcal{L}$ corresponds at computing the most likely sky $\vec{s}$ to correspond to the data $\vec{d}$ according to the noise $\vec{n}$. Equivalently, we can minimise thhe $\chi^2$ function associated with this likelihood : 

\begin{equation}
    \chi^2 \equiv -2\ln\mathcal{L} \propto (\vec{d} - H\vec{s})^{~T} N^{-1}(\vec{d} - H\vec{s}).
\end{equation}

To minimise this function, we need to compute the sky for which the derivate of $\chi^2$ function is zero. The minimization proceeds as follows: : 

\begin{align} \label{eq:mapmaking_inv}
    & \frac{\partial \chi^2}{\partial \vec{s}} = \frac{\partial \chi^2}{\partial \vec{s}^{~T}} = 0 \notag \\
     \iff  & \frac{\partial \chi^2}{\partial \vec{s}^{~T}} = \frac{\partial}{\partial \vec{s}^{~T}}(\vec{d}^{~T} N^{-1}\vec{d} - \vec{d}^{~T} N^{-1} H\vec{s} - \vec{s}^{~T} H^{~T} N^{-1}\vec{d} + \vec{s}^{~T} H^{~T} N^{-1}H\vec{s}) = 0 \notag \\
     \Rightarrow & - H^{~T} N^{-1} \vec{d} + H^{~T} N^{-1} H\vec{\hat{s}} = 0 \notag \\
     \Rightarrow & \vec{\hat{s}} = (H^{~T} N^{-1} H)^{-1} H^{~T} N^{-1} \vec{d},
\end{align}

where we note $\vec{\hat{s}}$ the maximal-likelihood solution of the map-making equation. 

Unfortunately, despite $\vec{\hat{s}}$ having an analytical solution, it requires to compute the inverse of the matrix $M \equiv (H^{~T} N^{-1} H)$. $M$ has a size $(N_{pix} x N_{pix})$, with $N_{pix} \sim 10^7$ for modern experiment. Inversion algorithms, like Cholesky factorisation or LU decomposition, have a computation cost in $\mathcal{O}(N_pix^{~3})$ and a memory usage in $\mathcal{O}(N_pix^{~2})$~\cite{golub13}. It then requires $\sim 10^{21}$ to invert such a matrix. With the CPU on my laptop\footnote{Intel i7-13850HX}, being able to perform around $100~GFLOP/s$, I would need about many years to perform an inversion ($\sim 320$ years with Cholesky decomposition). Even with exascale supercomputer, with $10^{18}~FLOP/s$, $\sim 20$ minutes are still needed. But, the main issue comes from memory, as in double precision\footnote{8 bytes per element}, the computation needs $8 \times (N_{pix})^2 = 800~TB$, just to store the matrix. \\
From these elements, matrix inversion is not feasible on realistic CMB datasets.

\subsection{Iterative Solver: PCG}

As discussed in the previous subsection, it is not reasonable to inverse the matrix $M$ to solve the map-making equation for realistic CMB datasets. Consequently, modern map-making algorithms rely on iterative solvers, such as the Preconditionate Conjugate Gradient (PCG) method, which avoid explicit matrix inversion and only require repeated evaluations of the operator. Indeed, equation \ref{eq:mapmaking_inv} can be rewritten : 

\begin{align} \label{eq:mapmaking}
    & \vec{\hat{s}} = (H^{~T} N^{-1} H)^{-1} H^{~T} N^{-1} \vec{d} \notag \\
    \Rightarrow & (H^{~T}N^{-1}H)\vec{\hat{s}} = H^{~T} N^{-1} \vec{d} \notag \\
    \iff & A\vec{x} = \vec{b},
\end{align}

with $A \equiv (H^{~T}N^{-1}H)$ and $b \equiv H^{~T} N^{-1} \vec{d}$. This linear equation is solvable using a Conjugate Gradient (CG), which will minimise the function $f(\vec{x}) = \frac{1}{2}\vec{x}^{~T}A\vec{x} - \vec{b}^{~T}\vec{x} + c$~\cite{liu2023}, with $c$ a constant. The gradient of $f$ is : 

\begin{equation}
    \Delta f(\vec{x}) = A\vec{x} - \vec{b}.
\end{equation}

CG algorithm will iterate on $\vec{x}$ to minimise $\Delta f(\vec{x})$, which exactly corresponds at finding the solution of equation \ref{eq:mapmaking}. Still, the size of $\vec{x}$ remains a problem, at the order of $N_{pix} \sim 10^7$. A preconditionning is then including in the Conjugate Gradiant solver. The preconditionner transforms the linear system in an equivalent one that converges faster, and can be seen as a matrix that approximates $A$, but easier to invert. A common preconditionner used in CMB experiment is $diag(A)$. A complete description of PCG algorithm can be found in \cite{shewchuk1994}.

\section{Numerical Implementation \& Methodology}

In the previous section \ref{section:mapmaking}, we derived the general solution of the map-making equation together with the iterative numerical solver required to compute it. We now extend this formalism to the specific case of the QUBIC instrument, with particular emphasis on its spectral imaging capabilities (see \ref{chap:spectral_imaging}). 

\subsection{Modelling \& Key Numerical Parameters}

In the previous derivation of the map-making equation, the frequency dependence of the instrumental model $H$ was neglected. For conventional imagers, this approximation is generally sufficient, since the instrumental response varies only weakly across the bandwidth. However, as discussed in the description of the spectral imaging capabilities of bolometric interferometry (in Chapter \ref{chap:spectral_imaging}), the frequency dependence plays a major role in shaping the QUBIC synthesized beam. The acquisition equation \ref{eq:acquisition} must therefore be written as : 

\begin{equation}
    \vec{d} = H\vec{s} + \vec{n} = \int_{\nu_{min}}^{\nu_{max}}\mathcal{H}_\nu\vec{s_\nu}d\nu + \vec{n}. 
\end{equation}

For the numerical modelisation, the matrix $\mathcal{H_\nu}$ cannot be integrated. Therefore, we can discretise the integration into frequency sub-bands, introducing the parameter $N_{sub}$, the number of sub-band : 

\begin{align} \label{eq:nsub_acq}
    \vec{d} &= \int_{\nu_{min}}^{\nu_{max}}\mathcal{H}_\nu\vec{s_\nu}d\nu + \vec{n} \notag \\
    & \approx \sum_{i = 1}^{N_{sub}} \mathcal{H}_{\nu_i} \vec{s_{\nu_i}}\Delta_{\nu_i} + \vec{n},
\end{align}
with $\Delta_{\nu_i}$ the bandwidth of the sub-bands associated with the frequency $\nu_i$.

To ensure the validity of such an approximation, $N_{sub}$ needs to be large enough. A deeper analysis of this parameter can be found in the subsection \ref{sub:nsub}. 
The discretisation in $N_{sub}$ sub-bands can be seen as a numerical trick to compute the integration over the frequency bandwidth. It is then necessary to compute the instrumental model at each frequency associated with the sub-bands, which is numerically expensive. $N_{sub}$ is then only a numerical, un-physical parameter. \\
Nevertheless, as we can separate the bandwidth into sub-bands, it is therefore possible to reconstruct them using the map-making solution. This idea, the fundamental basis of what we call the Frequency Map-Making, needs to reformulate the previous equation \ref{eq:nsub_acq} :

\begin{align}
    \vec{d} &= \sum_{i = 1}^{N_{sub}} \mathcal{H}_{\nu_i} \vec{s_{\nu_i}}\Delta_{\nu_i} + \vec{n} \notag \\
    &= \sum_{i = 1}^{N_{rec}} \left( \sum_{j = 1}^{f_{sub}} \mathcal{H}_{\nu_{ij}}  \Delta_{\nu_{ij}} \right) \vec{s}_{\nu_{i}} + \vec{n},
\end{align}

where $N_{rec}$ is the number of frequency maps we are trying to reconstruct and $f_{sub} \equiv N_{sub} / N_{rec}$, a sub parameter defining the number of frequency sub-bands integrated into each reconstructed maps.


\subsection{Angular resolution}

\subsection{External data ?}

We saw that Planck does not really improve FMM, while it is still useful to add at Xspectra stage.

\section{Characterisation}

\subsection{Numerical ressources: memory and parallelisation}

\subsection{Noise structure}

\subsection{Sub-bands Correlations}

\minitoc

\end{document}